\Needspace{12\baselineskip}
\subsubsection{3.2.15 UC15 安排并维护每日班次}\label{uc-v3-15}
\begin{BSRUsecase}{UC15 安排并维护每日班次}
主执行者 & 门店经理 \\
辅执行者 & 无独立外部执行者；数据库、SSE 与计时模块均为系统内部机制。 \\
涉众利益 & 经理需要覆盖门店各岗位；工作人员需要唯一且清楚的计划班次；组织需要保护已发生的出勤判断依据。 \\
目标与触发 & 经理选择日期或人员后创建班次，或编辑尚未产生打卡的班次。 \\
前置条件 & 已登录为经理；人员是非顾客角色；门店存在；日期和时间有效。编辑时原班次存在。 \\
成功后置条件 & 班次持久化且一人同日唯一；修改有审计；每日排班及个人班次经同步可见。 \\
最小保证 & 失败不提交部分业务变更；权限或版本冲突返回明确原因。业务事实以服务器记录为准。 \\
基本流程 & \begin{minipage}[t]{\linewidth}\raggedright 1 经理按工作日期和门店查看每日安排。\\ 2 系统显示人员、岗位、班次及当天打卡情况。\\ 3 经理选择人员、门店、日期与开始结束时间。\\ 4 系统校验 06:00–23:00 范围、先后关系和同日唯一性。\\ 5 编辑时系统同时校验原班次人员日期和新目标人员日期均无打卡。\\ 6 系统保存班次并记录操作者和调整内容。\\ 7 系统推送变更；经理每日安排与相关人员页面刷新。\strut\end{minipage} \\
扩展流程 & \begin{minipage}[t]{\linewidth}\raggedright 4a 同日已有班次：返回冲突，提示编辑原班次。\\ 4b 日期不存在、时间格式错误、开始不早于结束或越界：拒绝。\\ 5a 任一相关人员日期已有打卡：拒绝修改，不允许移动已有出勤的班次。\\ 6a 并发创建冲突：唯一约束使同一人员日期最多一条成功。\strut\end{minipage} \\
字段与状态事件 & id、userId、storeId、workDate、startTime、endTime；状态事件为班次创建、班次保存、冲突拒绝；具体岗位来自人员资料与账户角色。 \\
\rowcolor{GoldPale}业务规则 & BR-031 至 BR-034；排班面向全部非顾客工作人员，不将分包商自动转为雇员；一人每天单班，跨午夜和多段班次不在本轮。 \\
质量与验收 & 验证创建、编辑、日期边界、06:00 与 23:00、重复班次、已有打卡保护和非经理拒绝；经理每日表与个人班次数据一致。 \\
设计约束 & POST /api/shifts、PATCH /api/shifts/{id}；资料展示与员工新增、停用分开，当前不宣称已实现完整人事主数据 CRUD。 \\
\end{BSRUsecase}
